\section{A bounded-pattern estimate}
\label{sec:bounded-pattern}

All probabilities in this section refer to the distribution conditioned on
\(\operatorname{girth}(\Gamma)\ge L\).  We retain the predetermined vertices
\(x_A\) and \(x_B\). The possible rooted interval will accommodate a
null-homotopy of a root-to-vertex path in \Cref{sec:topology}.

Consider finitely many positive-length paths in \(\Gamma\), all cyclically
immersed and closed except possibly one immersed interval path whose
initial vertex is \(x_A\) or \(x_B\).  An \emph{occurrence} means a position
on one of these paths; different occurrences may traverse the same edge.
A pairing by \emph{interval pairs} consists of disjoint occurrence
intervals, paired in twos by length-preserving interval isometries.  Even
when the two intervals lie on the same path, they are distinct and
disjoint as occurrence intervals.  An order-preserving comparison
requires equal signed letters, and an order-reversing comparison
requires inverse letters in opposite orders.  We impose the additional
reduction condition
\begin{equation}
 \text{paired occurrences traverse distinct underlying undirected
 edges of }\Gamma.
 \label{eq:edge-reduction}
\end{equation}
Positions outside the paired intervals are called \emph{unpaired}.

\begin{samepage}
\begin{proposition}[Bounded-pattern estimate]
\label{prop:bounded-pattern}
There is a constant \(\varepsilon>0\), depending only on the fixed types
and turn weights, with the following property.  For every fixed positive
integer \(K\), fixed real \(C\ge1\), and fixed nonnegative integer \(I\),
the conditional probability tends to zero that there exists a system as
above having at most \(K\) paths, total length \(H\) with
\[
 L\le H\le CL,
\]
at most \(I\) interval pairs, and \(b\) unpaired occurrences, where
\(b\le\varepsilon H\).  In particular, \(\varepsilon\) is chosen before
\(K,C,I\), and these three bounds are held fixed as \(n\) tends to
infinity through its admissible values.
\end{proposition}
\end{samepage}

We first describe the actual image of a candidate system by boundedly
many chains. Sorting its edges by traversal multiplicity gives stages
whose total vertex exponent is nonpositive. Nearly aligned blocks turn
the path pairings into word constraints controlled by the squared-word
estimate. Together, these bounds make the product of the stage expectation
bounds exponentially small. Every realization must occur at every stage,
so one small expectation bound suffices for the union bound.

\begin{proof}
Fix \(K,C,I\).  Constants implicit in this proof may depend on these
bounds, whereas the final choice of \(\varepsilon\) will not.  Every
\(o(L)\) or \(o(s)\) estimate below is uniform over the patterns with these
fixed bounds.  Cutting comparison intervals at the chosen indexing
starts of closed paths changes only a fixed bound on the number of
interval pairs, so we make these cuts at the outset.

By \eqref{eq:reversal-weight}, reversing a chain does not change its word
weight.

\subsection*{The image graph and multiplicity stages}
Let \(\Delta\) be the actual image subgraph of a proposed system.  It has
at most \(CL\) edges and girth at least \(L\).  Mark every path start and
end, and every vertex whose degree in \(\Delta\) differs from two.  We
first show that \(\Delta\) consists of a bounded number of chains between
marks, with disjoint unmarked degree-two interiors.  A chain is allowed
to return to its initial mark.

Repeatedly delete leaves and their incident edges, and discard any edgeless
remnants. This preserves cycle rank. Set aside the components of the resulting core that are
circles; there are at most \(C\), since each uses at least \(L\) edges.
Suppress all degree-two vertices in the remaining core.  The resulting
multigraph has minimum degree at least three; its edges have positive
integer lengths, whose sum is at most \(CL\).  Suppose it has \(r>0\)
edges.  Start a non-backtracking walk uniformly on its \(2r\) directed
edges, and at every subsequent step choose uniformly among the allowed
continuations.  Uniform measure on directed edges is stationary: a
given directed edge with origin \(z\) has \(\deg(z)-1\) allowed
predecessors, each contributing
\(1/[2r(\deg(z)-1)]\).  This calculation also applies to loops and
parallel edges, with degrees counted by half-edges.

The expected expanded length of a walk of \(k\) edges is therefore at
most \(kCL/r\).  At least half of the probability lies on walks of
expanded length at most \(2kCL/r\).  Each particular walk has
probability at most \((2r)^{-1}2^{-(k-1)}\), so there are at least
\(r2^{k-1}\) such short walks.  Choose \(k=O(\log(r+1))\) so that this
number exceeds \((2r)^2\).  There are at most \((2r)^2\) ordered pairs of
endpoint vertices.  Two distinct short walks consequently have the same
endpoints.  Their expansions are distinct reduced paths.  Their union
contains a cycle, since in a tree there is only one reduced path between
two specified vertices.  That cycle has length at most \(4kCL/r\).
Girth at least \(L\) implies \(r\le4Ck\), bounding \(r\) in terms of
\(C\) alone.  Thus the cycle rank of \(\Delta\) is bounded.

The graph \(\Delta\) has no isolated vertices and has at most two leaves.
Indeed, a cyclically immersed path cannot visit a leaf, and an immersed
interval path can visit one only at an endpoint.  If \(\beta\) is the
cycle rank and \(c_\Delta\) the number of components, the degree-sum
identity gives
\[
 \sum_{\deg(z)\ge3}(\deg(z)-2)
 =2\beta-2c_\Delta+\#\{z:\deg(z)=1\}.
\]
It follows that the number and total degrees of branching vertices are
bounded.  There are also only boundedly many path-start and path-end
marks.  Every component contains a path start.  Suppressing all remaining
unmarked degree-two vertices now proves the asserted bound on the
number of marked chains, including circle components.

Every path traverses whole marked chains.  Write \(m_e\) for the total
number of traversals of an undirected edge \(e\) by all the paths; it is
constant on each chain.  If one immersed path traverses the same
oriented edge twice, the gap between the two starting positions is at
least \(L\): the intervening segment is a nonempty reduced closed walk
and hence contains a cycle.  A path of length \(h\) thus traverses any
given oriented edge at most \(h/L+1\) times.  Summing over paths and the
two orientations yields
\begin{equation}
 m_e\le 2(H/L+K)\le2(C+K).
 \label{eq:entropy-multiplicity}
\end{equation}
Consequently the total number of chain traversals is bounded as well.

There are only polynomially many unlabelled patterns in \(L\).  To
specify one, give the bounded marked multigraph, the positive integer
length of each chain, the side \(A\) or \(B\) of each component, the
oriented chain sequence of each path, the designation and starting side
of the possible interval path, and the comparison endpoints and their
orientation data.  Each integer endpoint has \(O(L)\) choices and there
are boundedly many of them.  We retain only patterns satisfying
\eqref{eq:edge-reduction}.  Vertex names in \(A\) and \(B\), and signed
edge labels, are not included in this description.  Crucially, the
abstract graph so described is the actual image graph: its subsequent
realizations must be injective on vertices and edges.

\paragraph{Multiplicity stages and vertex exponents.}
The density-counting viewpoint of organizing constraints by multiplicity,
discussed by Ollivier \cite{Ollivier2005}, motivates the stage organization
and selection below.  The conditioned matching estimate for the dependent
graph paths in this model is developed locally.

For \(1\le j\le m_*:=\max_e m_e\), let \(\Delta_j\) consist of all edges
with \(m_e\ge j\), together with their incident vertices.  Denote its
edge and vertex counts by \(E_j\) and \(V_j\).  Let \(\iota\) be one if
the interval path is present and zero otherwise.  Then
\begin{equation}
 \sum_{j=1}^{m_*}E_j=H,\qquad
 \sum_{j=1}^{m_*}V_j
 =\sum_{z\in\Delta^{(0)}}\max_{e\ni z}m_e
 \le H+\iota.
 \label{eq:entropy-incidences}
\end{equation}
For the inequality, let \(a(z)\) count interval-path endpoints at \(z\),
with multiplicity.  At a vertex, all traversal incidences except these
endpoints are coupled into passages between distinct half-edges, by
immersion and cyclic immersion.  Thus the incidences on a half-edge
of maximum multiplicity must be coupled to incidences on other
half-edges, except for at most \(a(z)\) endpoints.  Hence
\[
 2\max_{e\ni z}m_e
 \le \sum_{\text{half-edges at }z}m_e+a(z).
\]
Summing this inequality proves \eqref{eq:entropy-incidences}, because the
half-edge sum is \(2H\) and \(\sum_z a(z)=2\iota\).

Set \(k_1=\iota\), and \(k_j=0\) for \(j>1\).  In stage one the initial
vertex of the interval path has its prescribed image, saving one factor
of \(n\).  We claim no such saving at later stages.  The exponent needed
below satisfies
\begin{equation}
 \sum_{j=1}^{m_*}(V_j-E_j-k_j)\le0.
 \label{eq:entropy-n-exponent}
\end{equation}
In the later realization count, the exponent \(V_j-E_j-k_j\) will
count vertex-image choices against edge-prescription costs.
Equation~\eqref{eq:entropy-n-exponent} therefore controls the total
power of \(n\) across the multiplicity stages.
The remaining task is to obtain exponential decay from the label
constraints; counting only distinct edges would lose the incidence
balance that makes these two estimates compatible.

\subsection*{Word constraints on aligned blocks}
The aim is a common subdivision for which each comparison carries almost
all positions of a block into a single partner block. Orient each marked chain of length \(\ell\), identify it with
\([0,\ell]\), and put its edge positions at half-integers.  Splitting
comparisons at chain-traversal boundaries gives boundedly many
subcomparisons of the form
\[
 x\longmapsto\sigma x+a,\qquad \sigma\in\{1,-1\},\quad a\in\mathbb Z.
\]
These maps compare particular occurrences of subchains.  Relative to
the chosen chain orientations, their label relations are equality when
\(\sigma=1\) and inverse equality when \(\sigma=-1\).

Pad the list of integer offsets by zeros to a fixed length \(d\ge1\).
Put
\[
 N=\left\lfloor L^{1/(2(d+1))}\right\rfloor,
 \qquad R=\lfloor\sqrt N\rfloor.
\]
The simultaneous pigeonhole principle gives an integer
\(1\le u\le N^d\) such that
\(\operatorname{dist}(ua/L,\mathbb Z)\le1/N\) for every offset \(a\).
For completeness, place the \(N^d+1\) vectors of fractional parts of
\(k(a_1/L,\ldots,a_d/L)\), for \(0\le k\le N^d\), in the \(N^d\)
cubes of side \(1/N\).  Subtracting two vectors in the same cube gives
the required \(u\).  Choose the least such \(u\), so all choices are
determined by the unlabelled pattern, and set \(s=L/(uR)\).  Uniformly in
that pattern,
\begin{equation}
 \frac{L}{N^{d+1/2}}\le s\le\frac{L}{R},\qquad
 \max_a\operatorname{dist}(a/s,\mathbb Z)\le\frac{R}{N}.
 \label{eq:entropy-grid}
\end{equation}
In particular, \(s\) tends to infinity at least as a fixed positive
power of \(L\), while \(s/L\) and all the distances in
\eqref{eq:entropy-grid} tend to zero.

On each chain use the full grid intervals \([ks,(k+1)s)\) contained in
\([0,\ell]\), for integers \(k\ge0\).  The half-integer edge positions
in such an interval form a \emph{full block}.  Index these blocks over
all of \(\Delta\) by \(i\).  Their lengths and multiplicities satisfy
\[
 h_i=s+O(1),\qquad m_i=m_e\quad\text{on their chains}.
\]
The remaining end fragment of a chain has \(O(s+1)\) positions.  Every
traversal of the same chain uses this same subdivision, also when read
backwards.

Here is the alignment consequence of \eqref{eq:entropy-grid}, including
reflections.  Write \(a=\nu s+r\), with \(\nu\in\mathbb Z\) and
\(|r|=o(s)\).  A translation sends the \(k\)-th real grid interval to
the \((k+\nu)\)-th interval shifted by \(r\).  A reflection sends it to
\[
 ((\nu-k-1)s+r,(\nu-k)s+r].
\]
Because \(a\) is integral, the isometry preserves the half-integer
lattice.  In either case its image overlaps the indicated partner
grid block in a contiguous string of at least \(s-|r|-2=s-o(s)\)
positions.  The reversed endpoint convention in the reflection costs
only a bounded number of positions.

Figure~\ref{fig:block-alignment} illustrates the translation case.
The grid belongs to the underlying chain, not to an individual path
traversal: repeated traversals use the same blocks but give distinct
block occurrences. This distinction will matter when occurrence
pairings are projected to links between underlying blocks.

\begin{figure}[htbp]
\centering
\begin{tikzpicture}[x=0.85cm,y=0.85cm,>=Latex,font=\small,line width=0.6pt]
\fill[blue!12] (3,1.8) rectangle (6,2.12);
\draw[->] (0,1.96)--(9.3,1.96);
\foreach \position in {0,3,6,9}
  \draw (\position,1.7)--(\position,2.22);
\node[anchor=south] at (3,2.25) {$ks$};
\node[anchor=south] at (6,2.25) {$(k+1)s$};
\node[anchor=east] at (-0.2,1.96) {source};
\draw[->,dashed] (4.5,1.65)--(4.9,0.96);
\node[anchor=west] at (5.15,1.3) {$x\mapsto x+\nu s+r$};
\fill[orange!18] (3.4,0.65) rectangle (6.4,0.95);
\draw[orange!70!black] (3.4,0.65) rectangle (6.4,0.95);
\node[anchor=east] at (3.2,0.8) {translated block};
\fill[blue!12] (3,-0.14) rectangle (6,0.18);
\fill[blue!35] (3.4,-0.14) rectangle (6,0.18);
\draw[->] (0,0.02)--(9.3,0.02);
\foreach \position in {0,3,6,9}
  \draw (\position,-0.24)--(\position,0.28);
\draw[dashed] (3.4,-0.14)--(3.4,0.65);
\draw[dashed] (6.4,0.18)--(6.4,0.65);
\node[anchor=east] at (-0.2,0.02) {target};
\node[anchor=north] at (3,-0.28) {$(k+\nu)s$};
\node[anchor=north] at (6,-0.28) {$(k+\nu+1)s$};
\draw[<->] (3.4,-1.05)--(6,-1.05);
\node[anchor=north] at (4.7,-1.12) {overlap $s-r$};
\end{tikzpicture}

\caption{Approximate alignment under a translation, shown schematically
with \(0<r=o(s)\). A source block is translated by \(a=\nu s+r\);
its image overlaps the indicated target block in real length \(s-r\).
Restricting to half-integer edge positions loses at most a bounded
number of positions, giving \(s-o(s)\) compared letters. The shift is
exaggerated for visibility. The proof also treats negative shifts and
reflections.}
\label{fig:block-alignment}
\end{figure}

Join two full-block occurrences whenever one of the subcomparisons
identifies contiguous subwords in them of length at least
\((1-\kappa)s\), where \(\kappa=o(1)\) is chosen uniformly large enough
to absorb the alignment and rounding errors.  Record one such substring
comparison with each join.  These joins form a partial matching of
block occurrences: each compared substring occupies more than half its
block, whereas the original comparison domains are disjoint as
occurrence intervals.  Thus a block occurrence cannot participate in
two joins.

Let \(b_*\) be the number of unmatched full-block occurrences.  Only
boundedly many blocks lie within \(3s+O(1)\) of a subcomparison cut or a
chain endpoint.  Away from these cuts, a full block in a comparison
domain has its partner full block and hence is matched.  Every remaining
unmatched block away from the cuts consists entirely of unpaired
positions.  There are \(O(L/s)\) block occurrences in total, so rounding
their lengths and charging the exceptional blocks gives
\begin{equation}
 sb_*\le b+O(s+L/s)=b+o(L).
 \label{eq:entropy-unmatched-blocks}
\end{equation}

\paragraph{Links between underlying blocks.}
Forget occurrence indices in this matching, obtaining a multigraph
whose vertices are the full blocks \(i\).  We call its edges
\emph{links} and include isolated blocks.  The substring positions and
directions attached to every link are already fixed by the unlabelled
pattern.  In a component \(\mathcal U\) of this link graph, a word on any
one block determines the word on each next block along a spanning tree
except for \(o(s)\) letters.  Indeed, the recorded comparison specifies
\(s-o(s)\) letters of the next word, possibly with reversal and
inversion.  Since the alphabet is fixed, the number of extensions of a
root word to all blocks in \(\mathcal U\) is at most
\begin{equation}
 \exp\bigl(o(s)|\mathcal U|\bigr).
 \label{eq:entropy-propagation}
\end{equation}
Additional links and immersion constraints can only reduce this number.

If \(\mathcal U\) has a self-link, choose a self-linked block as the
root. Its word has only \(\exp(o(s))\) possibilities, or the pattern is
infeasible. To prove this, first consider a translation self-link. Its
constraints have the form
\[
 W(x)=W(x+t)\quad(x\in D),
\]
where \(D\) and \(D+t\) are contiguous subsets of the same block and
\(|D|=s-o(s)\).  The integral shift \(t\) is nonzero by
\eqref{eq:edge-reduction}; a zero shift would compare the same underlying
edge with itself.  On the \(h_i\) positions of the block, the
\(|D|\) distinct constraint edges form a forest of arithmetic chains
with step \(|t|\).  Thus this forest has \(h_i-|D|=o(s)\) components,
leaving at most \(|T|^{h_i-|D|}=\exp(o(s))\) root words.

For a reflection self-link, write its map as \(\phi(x)=a-x\).  The
intersection \(J=D\cap\phi(D)\) of its nearly full domain and image is
a contiguous reflection-invariant set with
\(|J|\ge2|D|-h_i>0\).  If \(a\) is odd, its center is a half-integer
position in \(J\), forcing that letter to equal its inverse.  If \(a\)
is even, the two half-integer positions on either side of its center
are adjacent positions in \(J\) and their letters must be inverse.
The first alternative contradicts the fixed-point-free involution on
\(T\), and the second contradicts immersion.  Thus reflection self-links
are infeasible.  Combining the translation bound with
\eqref{eq:entropy-propagation}, a component with a self-link has at most
\(\exp(o(s)|\mathcal U|)\) full word assignments altogether.

For any link component put
\[
 S_{\mathcal U}=\sum_{i\in\mathcal U}m_i,\qquad
 z_{\mathcal U}=\max_{i\in\mathcal U}m_i,\qquad
 b_{\mathcal U}=\#\{\text{unmatched block occurrences in }\mathcal U\}.
\]
If \(\mathcal U\) has no self-link, then
\begin{equation}
 2z_{\mathcal U}\le S_{\mathcal U}+b_{\mathcal U}.
 \label{eq:entropy-link-balance}
\end{equation}
Indeed, take a block of multiplicity \(z_{\mathcal U}\).  At least
\(z_{\mathcal U}-b_{\mathcal U}\) of its occurrences are matched, and
their partners are distinct occurrences of other blocks in the same
component.  There are only \(S_{\mathcal U}-z_{\mathcal U}\) such other
occurrences.  This proves \eqref{eq:entropy-link-balance}, also for
isolated blocks and in the presence of multiple links.

\paragraph{Binning the word weights.}
For each full-block word \(W_i\), read in its chain orientation, specify
the integer
\[
 d_i=\lfloor-\log P(W_i)\rfloor.
\]
We need only consider positive weights, since feasible block words are
immersed.  By Lemma~\ref{lem:squared-words}, every such word of length
\(h\) has \(P(W)\le e^{-\delta h}\) for all sufficiently large \(h\).
The smallest positive turn weight is a fixed positive number.  Hence
there is a constant \(a_0>0\), depending only on the turn weights, such
that for all large \(L\),
\begin{equation}
 \frac{\delta s}{2}\le d_i\le a_0s,
 \qquad
 \#\{W_i:\lfloor-\log P(W_i)\rfloor=d_i\}
 \le\exp(2d_i-\delta s).
 \label{eq:entropy-bins}
\end{equation}
For the counting assertion, every word in the indicated bin contributes
more than \(e^{-2d_i-2}\) to the sum of squared weights.  The same lemma
therefore bounds the number of these words by
\[
 \exp(2d_i+2-2\delta h_i)\le\exp(2d_i-\delta s),
\]
using \(h_i=s+O(1)\).  This also explains why the bin estimate uses the
squared, rather than the unsquared, word weights.

Enlarge each unlabelled pattern by its bin integers.  There are
\(O(L/s)\) blocks and \(O(s)\) possible bin values per block, so the
number of choices is at most
\[
 \exp\bigl(O((L/s)\log(s+1))\bigr)=\exp(o(L)).
\]
Together with the earlier polynomial count, the number of enlarged
patterns is \(\exp(o(L))\), uniformly for the fixed bounds.

\subsection*{Expected counts and stage selection}
The block estimates now control feasible strings. We combine them with
the matching-prescription estimate to count injective realizations of
each multiplicity stage, keeping string extension distinct from graph
embedding extension.

Fix an enlarged pattern.  A permitted label assignment on \(\Delta_j\)
is the restriction of some label assignment on all of \(\Delta\)
satisfying the bins, the immersion condition into the rose, and all
original comparisons.  This is a condition on strings: no embedding
of the inactive part of \(\Delta\) in the sampled graph is required.
Let \(X_j\) count injective realizations of \(\Delta_j\), with these
permitted labels and prescribed sides, and with the prescribed initial
vertex in stage one if \(\iota=1\).  A realization of the entire path
system gives \(X_j\ge1\) at every stage.

For one fixed permitted stage label assignment, its expected number of
vertex realizations is at most
\begin{equation}
 n^{V_j-E_j-k_j}
 \exp\left(-\sum_{i:m_i\ge j}d_i+o(L)\right).
 \label{eq:entropy-fixed-label-expectation}
\end{equation}
Here are the vertex and edge counts giving this bound.  A chain is
either entirely present or absent at stage \(j\), because its
multiplicity is constant.  Let \(Q_j\) be the number of active chains.
There are \(E_j-Q_j\) internal chain vertices.  If such a vertex has
incoming letter \(t\) and outgoing letter \(u\), then on its prescribed
side it has at most \(np\,w(t,u)\) possible images.  All other unfixed
vertices have at most \(n\) choices.  The fixed initial vertex, when
imposed, is a marked vertex.  Ignoring injectivity gives the upper bound
\[
 n^{V_j-k_j}p^{E_j-Q_j}
 \prod_{\text{internal chain turns}}w(t,u)
\]
on the number of vertex assignments.  We apply probability estimates
only to the feasible injective choices among these assignments.

Such a choice requires \(E_j\) distinct matching prescriptions.  Two
distinct edges cannot prescribe the same matching pair in either
direction, because injectivity and the immersion condition would then
identify their outgoing letters at an endpoint.  By
Lemma~\ref{lem:prescriptions}, the probability of all these prescriptions
is at most
\[
 \left(\frac{1+o(1)}{np}\right)^{E_j},
\]
uniformly since \(E_j=O(L)\).  Multiplication leaves the bounded factor
\(p^{-Q_j}\), the power \(n^{V_j-E_j-k_j}\), and the product of internal
turn weights.  Discard turn weights outside full blocks, since they
are at most one.  The retained product is
\(\prod_{i:m_i\ge j}P(W_i)\le\exp(-\sum_{i:m_i\ge j}d_i)\).
The bounded factor and the prescription errors are absorbed in
\(\exp(o(L))\), proving \eqref{eq:entropy-fixed-label-expectation}.

The number of permitted stage label assignments is at most
\begin{equation}
 \exp\left(o(L)+
 \sum_{\substack{\mathcal U\text{ without a self-link}\\
                         z_{\mathcal U}\ge j}}
       \left(2\min_{i\in\mathcal U}d_i-\delta s\right)\right).
 \label{eq:entropy-stage-labels}
\end{equation}
To see this carefully, consider a link component with at least one
active block, equivalently \(z_{\mathcal U}\ge j\).  The allowed labels
on its active blocks are projections of full feasible word assignments
on that component.  Their number cannot exceed the number of full
assignments.  If the component has no self-link, choose a root attaining
\(\min_{i\in\mathcal U}d_i\), count its words by
\eqref{eq:entropy-bins}, and extend along a spanning tree by
\eqref{eq:entropy-propagation}.  This argument is valid even if the
chosen root is inactive at this stage: it counts full word assignments
before projecting, and asserts no extension of an embedded graph.
A self-linked component has only its propagation error, as proved
above.  Components with no active block contribute nothing.  The sum
of propagation errors is \(o(L)\), since the total number of blocks is
\(O(L/s)\).  Arbitrary letters in the end fragments contribute at most
\(\exp(O(s))\) per marked chain, again \(\exp(o(L))\).  Ignoring any
further constraints between components gives
\eqref{eq:entropy-stage-labels}.

\paragraph{Selecting a stage and summing over patterns.}
Let \(B_j\) be the product of the right-hand sides of
\eqref{eq:entropy-fixed-label-expectation} and
\eqref{eq:entropy-stage-labels}.  Thus \(\operatorname{E}X_j\le B_j\).
Multiply these numerical upper bounds over
\(1\le j\le m_*\).  The total exponent of \(n\) is nonpositive by
\eqref{eq:entropy-n-exponent}.  For a link component without a self-link,
put \(d_{\min}=\min_{i\in\mathcal U}d_i\).  Its remaining contribution
to the logarithm of the product is
\[
 F_{\mathcal U}
 =-\sum_{i\in\mathcal U}m_i d_i
       +2z_{\mathcal U}d_{\min}-\delta s z_{\mathcal U}.
\]
If \(2z_{\mathcal U}\le S_{\mathcal U}\), then
\eqref{eq:entropy-bins} gives
\[
 F_{\mathcal U}
 \le-(S_{\mathcal U}-2z_{\mathcal U})\frac{\delta s}{2}
       -\delta s z_{\mathcal U}
 =-\frac{\delta}{2}sS_{\mathcal U}.
\]
If \(2z_{\mathcal U}>S_{\mathcal U}\), use
\eqref{eq:entropy-link-balance}, the upper bin bound, and
\(z_{\mathcal U}>S_{\mathcal U}/2\) to obtain
\[
 F_{\mathcal U}
 \le(2z_{\mathcal U}-S_{\mathcal U})a_0s
       -\delta s z_{\mathcal U}
 \le a_0s b_{\mathcal U}-\frac{\delta}{2}sS_{\mathcal U}.
\]
For a self-linked component the negative cost
\(-\sum_i m_i d_i\) alone is at most
\(-\delta sS_{\mathcal U}/2\).

Bounded multiplicities and boundedly many chains imply
\[
 \sum_{\mathcal U}sS_{\mathcal U}=H+O(L/s+s)=H+o(L),
 \qquad
 s\sum_{\mathcal U}b_{\mathcal U}\le b+o(L),
\]
where the second inequality is \eqref{eq:entropy-unmatched-blocks}.
There are only boundedly many stages by
\eqref{eq:entropy-multiplicity}, so all their error terms remain
\(o(L)\).  We have proved
\begin{equation}
 \log\prod_{j=1}^{m_*}B_j
 \le-\frac{\delta}{2}H+a_0b+o(L).
 \label{eq:entropy-product-bound}
\end{equation}

Choose, once and for all,
\[
 \varepsilon=\min\left\{\frac14,
              \frac{\delta}{16(1+a_0)}\right\}>0.
\]
This depends only on the fixed turn weights.  For \(b\le\varepsilon H\)
and all sufficiently large \(L\), uniformly over the fixed pattern
bounds, \eqref{eq:entropy-product-bound} gives
\[
 \prod_{j=1}^{m_*}B_j\le e^{-\delta H/4}.
\]
Put \(M=\max\{1,\lceil2(C+K)\rceil\}\), so \(m_*\le M\).  Some stage
therefore has \(B_j\le e^{-\delta H/(4M)}\).  Since a realization of
the entire enlarged pattern requires \(X_j\ge1\) at every stage,
Markov's inequality gives
\[
 \operatorname{Pr}(\text{this enlarged pattern is realized})
 \le\min_j\operatorname{E}X_j
 \le\min_j B_j
 \le e^{-\delta H/(4M)}.
\]
This uses no independence between stages.  We multiplied numerical
bounds solely to find one stage with a small expectation.

Finally there are \(\exp(o(L))\) enlarged patterns and \(H\ge L\).
The union bound is at most
\(\exp(o(L)-\delta L/(4M))\), which tends to zero.  This proves the
proposition with the asserted order of quantifiers.
\end{proof}
